\section{Ordinary relative Iitaka reductions}\label{ordinary:reductions}

The absolute construction proves the main theorem directly. For
projective fibrations it is also useful to start with the relative
Iitaka fibration and retain all directions of a second Iitaka base.
We derive that formulation from Theorem~\ref{num:main} and
identify its parameter field inside the original $\C(Y)$. The
additional companion inputs in this section are the relative
rank-one section comparison and its reduced-boundary cyclic refinement;
their exact outputs are specified in Proposition~\ref{input:reduction}.
The reduced boundary is essential to its entire-top-space assertion.
The numerical construction itself used an arbitrary nef divisor;
here the relative Kodaira-zero comparison identifies that divisor
with a parabolic Hodge line. At the end, we identify the resulting
fields inside $\C(X)$, so the conclusion concerns the original
whole geometric generic fiber.

\begin{proposition}[Ordinary relative reduction]\label{input:reduction}
Let $f:X\to Y$ be a surjective morphism with connected fibers between
smooth connected projective complex varieties. Assume $\kappa(Y)\geq0$,
and put $d=\kappa(F)\geq0$ for the geometric generic fiber $F$. On smooth projective models there is a
factorization
\[
 X'\xrightarrow{x}W\xrightarrow{g}Y
\]
with connected geometric generic fibers and $\dim(W/Y)=d$.
There are an effective SNC rational boundary $B$ on $W$, a nef
rational divisor $M$, and a morphism $h_0:W\to Z$ with connected
fibers, with $Z$ smooth projective, such that
\begin{enumerate}[label=\textup{(\roman*)}]
\item $K_W+B+M$ is big over $Y$;
\item $M\sim_{\Q}h_0^*A$ for a nef rational divisor $A$ on $Z$,
      and $K_Z+aA$ is big for every sufficiently large rational $a$;
\item the divisible section systems of $K_W+B+M$ identify with
      those of $K_X$, preserving their rational-map dimensions;
\item the generic pair $(J,D_J)$ of $x$, with the reduced exceptional
      boundary inherited from $X'\to X$, has logarithmic Kodaira
      dimension zero. Its minimal-index cyclic root cover is
      geometrically connected and has entire logarithmic top-form
      space of dimension one, with every positive logarithmic
      plurigenus at most one;
\item $M$ is the rational parabolic extension of this entire top
      line, identified with the top line of the unique rational
      polarized pure weight grade carrying it.
\end{enumerate}
On an initial reduction model $B$ is the normalized boundary of
\cite[\PoneNormalization]{SubadditivityCompanion}. Subsequent source models of
$W$ carry its strict transform plus reduced exceptional boundary,
and the pullback of $M$; their section systems remain those of the
original $X$.
\end{proposition}

\begin{proof}
Take the relative Iitaka fibration of $K_X$ over $Y$, using the
relative algebraic closure of its section field, and resolve its
maps. The Iitaka fibration theorem gives relative dimension $d$
and Kodaira dimension zero on its geometric generic fiber
\cite[Section~6]{IitakaOriginal}. Give $X'$ the strict transform of
the zero boundary plus the reduced exceptional divisor. The
multiplicative section-ring identity in
Lemma~\ref{setup:log-modification}, also relative over $Y$, identifies
its relative log Iitaka fibration with the one just chosen.
Consequently $(J,D_J)$ has logarithmic Kodaira dimension zero.

Apply the rank-one comparison of \cite[\PoneSectionComparison]{SubadditivityCompanion}, with the original $X$
as its reference model. The relative section comparison gives
$K_W+B+M$ the full section field of dimension $d$ on the generic
fiber over $Y$; hence it is relatively big. It also gives (iii).
The normalization lemma in \cite[\PoneNormalization]{SubadditivityCompanion}
identifies its initial boundary using two different order tests.
The coefficient defining sections is the minimum over the actual
components of a fiber, whereas the discriminant threshold also
allows divisorial valuations on higher models. The stated
normalization compares these quantities on the same SNC model;
we do not replace one by the other. The reduced-boundary results \cite[\PoneReducedRoot]{SubadditivityCompanion} and
\cite[\PoneReducedModuli]{SubadditivityCompanion} give (iv)--(v), including
the normalization of $M$ without an extra cyclic-index factor.

The entire highest line in (v), with its actual parabolic
extension, satisfies the hypotheses of Proposition~\ref{hd:adjoint}.
That proposition gives $h_0$, $Z$ and $A$ after smooth projective
modifications. Its positive value of $a$ suffices for all
larger $a$, since $A$ is nef. Use
Lemma~\ref{prelim:sections}, with the pulled-back rational twist, for the boundary on the modified
$W$. Its section comparison, and the original-reference convention
in \cite[\PoneSectionComparison]{SubadditivityCompanion}, preserve (i) and (iii).
The parabolic line pulls back functorially, so (v) is preserved as
well. The modifications and the relative Iitaka construction do not
change geometric generic integrality.
\end{proof}

\begin{corollary}[The ordinary second Iitaka construction]
\label{input:second-iitaka}
For the data of Proposition~\ref{input:reduction}, there are
a morphism $q:W\to U$ with connected fibers to a smooth projective
variety $U$ and a morphism
$j:U\to Z$ with $h_0=j\circ q$. The general $q$-fiber $G$ satisfies
\[
 \kappa(G,K_G+B|_G)=0,
 \qquad G\longrightarrow g(G)\text{ is generically finite}.
\]
There are smooth projective varieties $T_0,S$ and a diagram
\[
\begin{tikzcd}[column sep=large]
 W\arrow[r,"{(g_S,q)}"]\arrow[d,"q"'] &
 (S\times_{T_0}U)_{\mathrm{main}}\arrow[r]\arrow[d] &
 S\arrow[r,"\rho"]\arrow[d] &Y\\
 U\arrow[r,equal]&U\arrow[r]&T_0&
\end{tikzcd}
\]
with $g=\rho\circ g_S$, where $\rho$ is birational and the map
to the main product component is dominant and generically finite.
The geometric generic fibers of $S\to T_0$ and $U\to T_0$ are
integral, and
\begin{equation}\label{eq:imported-dimension}
 \kappa(S_t)=0,
 \qquad \kappa(X)=\kappa(W,K_W+B+M)=\dim U=d+\dim T_0.
\end{equation}
Here $t$ is a very general point of $T_0$.
All models use the section-preserving boundary convention of
Proposition~\ref{input:reduction}.
\end{corollary}

\begin{proof}
Take $M=h_0^*A$ in Theorem~\ref{num:main}; rational linear
equivalence suffices there. Its conclusions, with its
product diagram~\eqref{num:product} and dimension
identity~\eqref{num:dimension-conclusion}, give the asserted diagram
and dimension formula. Its hypotheses are exactly
(i)--(ii) of Proposition~\ref{input:reduction}, the effective SNC
boundary, and $\kappa(Y)\geq0$, with $D_Y=0$ in that theorem.
The section-system identification
in Proposition~\ref{input:reduction}(iii) transfers the resulting
dimension to $X$. This retains the actual image-cycle parameter
space, before its dimension is compared with $\kappa(Y)$.
\end{proof}

\begin{corollary}[Ordinary cyclic covers along the second Iitaka fibers]
\label{prop:ordinary-period}
In the notation of Proposition~\ref{input:reduction} and
Corollary~\ref{input:second-iitaka}, let $J$ be the generic fiber of
$x:X'\to W$ over $K=\C(W)$. The minimal logarithmic generator
$\omega\in H^0(J,p(K_J+D_J))$ is an ordinary regular pluriform.
Its resolved connected cyclic root cover $\widetilde J$ satisfies
\[
 H^0(\widetilde J,aK_{\widetilde J})=K\eta^a
 \qquad(a\geq1),
\]
where $\eta$ is its root top form. In particular,
$\kappa(\widetilde J)=0$ and $h^0(\widetilde J,K_{\widetilde J})=1$.

On a dense smooth open of every very general fiber of $q:W\to U$,
the family of these covers is geometrically birationally constant
with its deck action, and the same holds for the $J$-family after
taking the quotient. More precisely, over the geometric generic
$q$-fiber these families become birational to varieties over
$b_U=\overline{\C(U)}$ after a finite extension of that fiber's
function field, equivariantly for the deck group in the cover case.
These identifications persist under algebraically closed field extension.
\end{corollary}

\begin{proof}
Choose a nonzero ordinary pluricanonical form on the generic fiber
of $f$, which has Kodaira dimension $d\geq0$. Pull it to the induced
smooth birational model and restrict to the generic fiber of its
map to $W_y$. Adjunction, after dividing by a local determinant
frame on $W_y$, gives a nonzero ordinary form
$\theta\in H^0(J,mK_J)$ for some $m>0$. The restriction is nonzero
because a proper zero divisor cannot contain the generic fiber.
The logarithmic section space in degree $mp$ is one-dimensional,
so
\[
                  \theta^{\otimes p}=c\,\omega^{\otimes m}
                  \qquad(c\in K^*).
\]
Testing divisorial orders shows that $\omega$ has no poles. Its
index remains the original minimal $p$; the common degree is used
only to compare the forms.

In the reduced cyclic-cover lemma of \cite[\PoneReducedRoot]{SubadditivityCompanion}, this gives
\[
 E_J=\operatorname{div}(\omega)+pD_J,
 \qquad \Delta_J=D_J-E_J/p
       =-p^{-1}\operatorname{div}(\omega)\leq0.
\]
Thus $\Delta_J^{=1}$ and the specified pole boundary on the root
cover are empty. That lemma proves the displayed equality of
\emph{ordinary} pluricanonical spaces. Moreover,
\cite[\PoneReducedModuli]{SubadditivityCompanion} now identifies $M$
directly with the rational parabolic extension of the entire
compact top Hodge line. There is no open-fiber boundary to compare.

Spread an equivariant smooth projective model of the cover over a
nonempty smooth open $W^\circ\subset W$, shrinking it to remove
any vertical remnants of the spread pole boundary. Let $G$ be a very general
smooth $q$-fiber meeting this open. Since $h_0=j\circ q$ and
$M\sim_{\Q}h_0^*A$, the degree of $M$ on every complete test curve
in $G$ is zero. Lemma~\ref{hodge:extensions} and
Corollary~\ref{hodge:degree-zero} therefore make the projective
compact top line constant on each such curve. Complete-intersection
curves through a general point with tangent directions spanning
$T_G$ show that the line is constant on $G\cap W^\circ$.
Corollary~\ref{hodge:restricted-character}, applied anew on this
restricted locus, gives finite character.

Apply Corollary~\ref{cc:loci} to $W^\circ\to U$, using its
relative assertion and the finite deck action. Its hypotheses on
geometric generic integrality follow from
Corollary~\ref{input:second-iitaka}. This proves the stated
constancy of the covers, including over $b_U$. Taking invariant
function fields gives the assertion for $J$: invariants under a
finite group commute with field extension in characteristic zero.
\end{proof}

\begin{corollary}[The absolute Iitaka base and descent of the whole fiber]
\label{ordinary:absolute-bases}
\label{prop:full-descent}
The variety $U$ in Corollary~\ref{input:second-iitaka} is a
birational model of the absolute Iitaka base of $X$. Inside
$\C(X)$, the field $\C(W)$ is the relative algebraic closure of
$\C(Y)\C(U)$. Consequently the field $\C(T_0)\subset\C(Y)$
in that Corollary agrees with the field obtained from
Propositions~\ref{it:diagram} and~\ref{it:dimensions} applied to
the projective morphism $f$ with empty original boundaries.

In particular, inside the prescribed algebraic closure $\Omega$
of $\C(Y)$, put $b=\overline{\C(T_0)}$. There is a finitely
generated field $E_b/b$ such that
\begin{equation}\label{ordinary:whole-field}
 \operatorname{Frac}(\Omega\otimes_bE_b)
 \simeq
 \operatorname{Frac}(\Omega\otimes_{\C(Y)}\C(X)).
\end{equation}
Thus the whole geometric generic fiber of $f$ is birationally
defined over this specific field $b$, and
$\operatorname{Var}(f)\leq\dim T_0$.
\end{corollary}

\begin{proof}
We identify first the absolute section field, then the intermediate
field of the combined map, and finally the intersection of the
original base field with the absolute section field. All these
identifications take place inside $\C(X)$.

Put $L=K_W+B+M$, and let $E\subset\C(W)$ be its stabilized
section field. Since $L|_G=K_G+B_G$ has Iitaka dimension zero,
every section ratio is constant on the general $q$-fiber $G$.
Connectedness of $q$ gives $E\subseteq\C(U)$. The equality
$\kappa(W,L)=\dim U$ makes $\C(U)/E$ a finite extension.
The stabilized section field is relatively algebraically closed
in $\C(W)$, by the Iitaka theorem as used in
Proposition~\ref{it:diagram}. Hence $E=\C(U)$. The multiplicative
section comparison in Proposition~\ref{input:reduction} identifies
this field with the absolute canonical section field of $X$.

The product diagram in Corollary~\ref{input:second-iitaka}, with
$S\to Y$ birational, shows that $\C(W)$ is finite over
$\C(Y)\C(U)$. It is relatively algebraically closed in $\C(X)$
because $x$ has connected fibers. These two facts characterize
the relative algebraic closure used in
Proposition~\ref{it:diagram}, so the intermediate fields coincide.

It remains to identify the parameter field with its embedding,
rather than only its transcendence degree. In either construction,
the product diagram and geometric generic integrality give an
injection
\[
 \C(Y)\otimes_{\C(T_0)}\C(U)\lhook\joinrel\longrightarrow\C(W).
\]
For $u\in\C(Y)$ and $v\in\C(U)$, linear independence over
$\C(T_0)$ shows that $u\otimes1=1\otimes v$ forces
$u=v\in\C(T_0)$. Consequently
\[
                \C(T_0)=\C(Y)\cap\C(U)\quad\text{inside }\C(X).
\]
The two larger fields have already been identified, so the parameter
fields coincide with their embeddings in $\C(Y)$. Proposition~\ref{ds:whole}
now gives Equation~\eqref{ordinary:whole-field} over this exact
field $b$. The variation bound follows from its definition.
\end{proof}

\subsection{Comparing compact and open top lines}

Corollary~\ref{prop:ordinary-period} used the empty pole boundary
to identify the moduli line with the compact top line. The following
comparison treats the additional case in which a relative SNC
boundary is retained on a smooth cover. If the entire compact and
open top-form spaces are both one-dimensional, it transfers
flatness through the pure compact image and retains a prescribed
finite group action.

\begin{lemma}[Compact-to-open comparison]
\label{ordinary:compact-open}
Let $\pi:P\to T$ be a smooth projective morphism of relative
dimension $r$ with connected fibers, where $T$ is a smooth connected
complex algebraic variety. Let $D\subset P$ be a relative reduced
SNC divisor, with every stratum smooth over $T$, and suppose
\[
 h^0(P_t,K_{P_t})=h^0(P_t,K_{P_t}+D_t)=1
 \qquad(t\in T).
\]
Write $\mathbb H=R^r\pi_*\Q$, let $\mathbb V$ be the
degree-$r$ cohomology variation of $P\setminus D$, and put
$\mathbb Q=\im(\mathbb H\to\mathbb V)$.
Then $\mathbb Q=W_r\mathbb V$ is a polarizable pure variation,
and the compact-to-open map identifies its top line $F^r\mathbb Q$
with both $F^r\mathbb H$ and the entire top line $F^r\mathbb V$.
If $F^r\mathbb Q$ is flat, the compact top line $F^r\mathbb H$
is flat. A prescribed finite group acting on $(P,D)$ over $T$
admits an equivariant horizontal Hodge splitting realizing this
identification.
\end{lemma}

\begin{proof}
The logarithmic Hodge description identifies the map on top pieces
with the nonzero inclusion
\[
 H^0(P_t,K_{P_t})\lhook\joinrel\longrightarrow
 H^0(P_t,K_{P_t}+D_t);
\]
see \cite[Corollaries~3.2.13(ii) and~3.2.14]{Deligne71}.
The dimensions make it an isomorphism. Degree-$r$ open cohomology
has weights at least $r$, and its compact image is precisely
$W_r\mathbb V$ \cite[Corollary~3.2.17]{Deligne71}.
Strictness therefore identifies the three top lines. In particular,
the only weight grade carrying the logarithmic top line is this
pure weight-$r$ image.

The kernel of $\mathbb H\twoheadrightarrow\mathbb Q$ is a
pure Hodge subvariation. Its orthogonal complement for a
polarization of $\mathbb H$ is horizontal and is a Hodge
subvariation, and maps isomorphically to $\mathbb Q$; this is
the polarized splitting argument in \cite[Lemma~4.2.3(i)]{Deligne71}.
It supplies a horizontal Hodge right inverse $j$. Since the
kernel has zero $F^r$ piece, $F^r\mathbb H=j(F^r\mathbb Q)$.
Thus flatness of the latter line gives flatness of the former.

For a finite group $\Gamma$, replace $j$ by
\[
 j_\Gamma=\frac1{|\Gamma|}\sum_{\gamma\in\Gamma}
       \gamma_{\mathbb H}\,j\,\gamma_{\mathbb Q}^{-1}.
\]
Every summand is a horizontal Hodge right inverse, so their
average is one as well, and reindexing the sum proves equivariance.
The rank-one top identification is unchanged. This proves the
comparison with the finite action retained.
\end{proof}
